\chapter{The played board}

\textbf{Purpose:} Separate what the rules put on a board from what the players did, and report the
negative that came out. \textbf{Scope:}
\texttt{examples/\allowbreak{}game\_\allowbreak{}theory/\allowbreak{}\{1\_\allowbreak{}represent,3\_\allowbreak{}reference,4\_\allowbreak{}measure\}/\allowbreak{}}

\section{Why a board is worth reading at all}

Every corpus in this work holds two things at once and cannot separate them. A text carries its
grammar and its author together, and there is no way to obtain the same author writing under
different grammar. The workbook's Salishan rows spend real effort on exactly this and the effort does
not fully succeed.

A board separates them for nothing. The rules are fixed and public. The play is a choice. Three
players of measurably different strength on one rule set give nine corpora that differ in the second
alone, and the rule set's contribution is in all nine.

\section{The representation, and what happens to an empty square}

A square is a point and the piece on it is its value. Reversi on sixty four squares, two colors.

Emptiness does not get a symbol. A board part way through a game is mostly empty, and a symbol for
emptiness puts the game's clock into the histogram, where a measure reading counts will find it and
report how far along the game was as though it were arrangement. Measurements that read values drop
the empty squares; the one that reads the grid's shape keeps them, because a hole in a row is part
of the shape.

The null is this work's null permutation written for a grid: the same pieces on the same squares in
a uniformly drawn order. It keeps how many of each color are on the board and which squares are
occupied, and destroys only which color sits where. Everything a rule imposed is in what it
destroys.

The measure is the share of occupied squares whose neighbor one step away carries the same color,
averaged over both axes. The dispersion measure this work usually quotes needs four symbols clearing
an occurrence floor, and a board has two colors. This measure reads the same quantity at the only lag
a board is large enough to carry.

\section{The departure, and the control that says what it is}

Two hundred games with both sides drawing uniformly from the legal moves. Live 0.6555 against a null
of 0.5350, a ratio of 1.225. Six disjoint batches of two hundred give 1.2253, 1.2193, 1.2219, 1.2238,
1.2222 and 1.2282, a spread of 0.0089. The departure is about twenty five times the spread between
batches.

A departure on its own does not say which ingredient produced it, and a board supplies the control
for nothing. Play the same board with the same two colors and the same number of squares filled,
where a move places a piece and nothing is turned over. That game has rules and has play and has no
flip.

\begin{center}
\begin{tabular}{@{}lrrr@{}}
\toprule
& live & null & ratio \\
\midrule
Reversi, 200 games & 0.6555 & 0.5350 & 1.225 \\
place-only, 400 boards, from Reversi's opening & 0.4763 & 0.4904 & 0.971 \\
place-only, 400 boards, from an empty board & 0.4915 & 0.4923 & 0.998 \\
\bottomrule
\end{tabular}
\end{center}

The departure is the flip rule. A board being a board contributes nothing.

\subsection{The 0.971 row, which was a mistake of ours}

The first version of that control kept Reversi's four fixed center pieces and read 0.971, three
floors below the null where the argument wanted one. The gap is the opening itself. Those four
squares are laid out black over white and white over black. All four neighboring pairs inside them
therefore disagree by construction, where a draw would have half of them agree. Four guaranteed
disagreements out of a hundred and twelve pairs is 0.018 of the reading and accounts for all of the
gap. Starting from an empty board returns 0.998.

It is kept because it points the other way as well. The same four squares sit under every Reversi row
above and push them in the same direction. The ratio of 1.225 is understated by about that much and
is never overstated.

\section{The negative}

The three arms are the textbook ones. Drawing uniformly from the legal moves knows the rules and
has no strategy. Taking the move that turns over the most pieces is the standard weak heuristic and is
weak here too. Taking a corner where one is offered, and otherwise turning over the fewest pieces, is
the standard improvement: a corner can never be flipped back, and a small flip leaves the opponent
fewer replies.

Strength is not asserted. It is played out, and the win share is an independent quantity that is not
a reading of the departure.

\begin{center}
\begin{tabular}{@{}lrrr@{}}
\toprule
black v white & live & ratio & black's share of wins \\
\midrule
random v random & 0.6555 & 1.225 & 0.380 \\
random v greedy & 0.6581 & 1.191 & 0.398 \\
random v corner & 0.6722 & 1.224 & 0.320 \\
greedy v random & 0.6543 & 1.208 & 0.630 \\
greedy v greedy & 0.6404 & 1.184 & 0.492 \\
greedy v corner & 0.6735 & 1.198 & 0.425 \\
corner v random & 0.6700 & 1.241 & 0.645 \\
corner v greedy & 0.6881 & 1.194 & 0.672 \\
corner v corner & 0.6723 & 1.255 & 0.532 \\
\bottomrule
\end{tabular}
\end{center}

The win shares order the arms cleanly and in the expected direction. Corner beats random 0.645 and
beats greedy 0.672; greedy beats random 0.630; each arm against itself sits near a half.

The departures do not order them at all. They run 1.184 to 1.255 with no relation to the column
beside them, and the correlation over the nine pairings is 0.067.

The explanation is the same one that makes the flip rule the source of the departure. Every legal
move turns over a bracketed run, whoever chose it. A strong player and a weak one leave the same kind
of trace and differ only in where on the board they put it. The measure reads what was done to
the board. It does not read who chose it.

\section{Reading the width back off it}

The picture work claims an image read as a byte sequence returns its own width. A board is the same
shape of object with a width nobody has to look up.

\begin{center}
\begin{tabular}{@{}llrr@{}}
\toprule
arm & reading & period & margin \\
\midrule
random & row major & 8 & 0.0917 \\
random & column major & 8 & 0.1291 \\
random & scattered null & 7 & 0.0106 \\
corner & row major & 8 & 0.1209 \\
corner & column major & 8 & 0.1368 \\
corner & scattered null & 10 & 0.0093 \\
\bottomrule
\end{tabular}
\end{center}

Eight, across the rows and down the columns, on a weak arm and a strong one, at seven to fifteen
times what the same pieces scattered reach. The null returns a wrong number at a margin near zero.
That is the correct behavior: it is asked for a period and reports one with no confidence at all.

That result has a boundary and the next chapter is where it sits.
